\chapter{Conclusion}\label{ch:conclusion}
This thesis reads a recording as score, expression, timbre and noise, and
recovers the expression layer note by note together with its uncertainty.
The proposed model is a single multi-output Gaussian process on the score
graph, with the channels coupled by coregionalization, the side information
marginalized into the kernel, and every hyperparameter learned by the
per-piece evidence. It contains the strongest two-stage pipeline as a
special case, the diagonal-$B$, fixed-graph, plug-in-mean, no-feature slice
of its parameter space on the shared-shape restriction, and it recovers
hidden expressive detail better than that pipeline. This prediction was
fixed in advance and is supported on a held-out test set of pieces that no
modelling decision had touched.

The marginal value of each ingredient is isolated. Per-piece Bayesian
weighting of score features and music-model embeddings does the recovering,
and the graph makes the uncertainty calibrated, a contribution that the
held-out test set also supports. The music model earns its place as a
feature provider inside the kernel. One step further, on the development
set of both corpora, the graph's value is reproduced by the same graph
without its pitch term, so the structure rewarded is smoothness in score
time over about two beats. Six downstream tasks draw two boundaries:
structure helps note by note but not for whole-performance summaries, and
per-piece adaptation helps interpolation but not excerpt extrapolation. The
evaluation is protected against leakage, contamination, fit collapse and
selection inflation.

The piano contribution stands on its own, and the same prior has been
carried, unchanged in form, to intonation, vibrato, loudness, timing and
vibrato onset delay on real recordings. On the strings-and-winds held-out
test set the predictions for intonation recovery, vibrato calibration and
coverage are supported, and the secondary timing hypothesis is not. Its diagnosed remedy, correlated timing noise across
notes, was then supported on a second, disclosed evaluation of the same test
set. A tonal-metric calibration hypothesis evaluated alongside it was not
supported.

The step to audio through a differentiable synthesizer is where the most
novel work lies. Its development studies reach whole pitch curves whose
median error against the reference centre is below that of the tracker's
frames, and note boundaries inferred jointly from score-nominal starting
points. They also fix the design constraint that the remaining joint loop
must solve, bridging the estimand gap between the harmonic model and the
annotation chain. The score prior has entered that inference over three
spans, with a two-sided answer. A note or joint with no acoustic evidence
inherits its context's common offset with a calibrated band, and the score
graph outperforms a purely temporal prior at phrase scale, but nothing
measurable beyond the common offset has yet been found. The strong form of the hypothesis is left for the joint loop.

% =====================================================================
